% !TEX program = lualatex
% ECON 201, Worksheet 3: Making Economic Decisions. Compile with LuaLaTeX so
% the PDF is tagged (see syllabus/Econ201-Syllabus.tex for the same setup).
% Build from the course root:
%   latexmk -lualatex -cd worksheets/worksheet03.tex && latexmk -c -cd worksheets/worksheet03.tex
\DocumentMetadata{
  pdfstandard=UA-2,
  pdfversion=2.0,
  lang=en-US,
  tagging=on
}
\documentclass[11pt]{article}

\usepackage[letterpaper, margin=0.8in, top=0.6in, bottom=0.55in]{geometry}
\usepackage{fontspec}
\setmainfont{Fira Sans}
\usepackage{amsmath}
\usepackage{amssymb}
\usepackage{graphicx}
\usepackage{booktabs}
\usepackage{array}
% Questions are labelled by hand: under tagging=on, enumitem's enumerate
% did not step its counter, so every item printed the same label.
\newcounter{q}
\newcommand{\Q}[1]{\stepcounter{q}\par\needspace{5\baselineskip}\vspace{3pt}\noindent\hangindent=1.7em\hangafter=1
  \makebox[1.7em][l]{(\alph{q})}#1\par}
\usepackage{xcolor}
\usepackage{needspace}
\definecolor{accent}{HTML}{BF5700}
\usepackage{titlesec}
\titleformat{\section}{\large\bfseries\color{accent}}{}{0pt}{}
\titlespacing*{\section}{0pt}{6pt}{2pt}
\usepackage{hyperref}
\hypersetup{pdftitle={ECON 201 Worksheet: Making Economic Decisions}, pdfauthor={Div Bhagia}, hidelinks}
\pagestyle{empty}
\setlength{\parindent}{0pt}
\setlength{\parskip}{3pt}

% Ruled answer space: n lines. Plain TeX loop; pgffor's \foreach breaks the
% enumerate counter when used inside an item. Light grey rules, indented to
% the question text and spaced for handwriting, so they read as writing
% guides rather than section dividers.
\definecolor{answerrule}{HTML}{B8B8B8}
\newcount\answerlinecount
\newcommand{\answerlines}[1]{%
  \par\answerlinecount=#1
  \loop\ifnum\answerlinecount>0
    \vspace{16pt}\noindent\hspace*{1.7em}%
    {\color{answerrule}\rule{\dimexpr\linewidth-1.7em\relax}{0.4pt}}\par
    \advance\answerlinecount by -1
  \repeat}

% A cell to be filled in. Same width everywhere so the columns line up.
\newcommand{\blank}{\rule{1.15in}{0.3pt}}

\begin{document}

{\LARGE\bfseries\color{accent} Worksheet: Making Economic Decisions}\hfill{\large ECON 201}

\vspace{2pt}
Write your answers in the spaces.

\needspace{7\baselineskip}\section{Definitions from today}

\begingroup\small\setlength{\parskip}{2pt}
\textbf{\color{accent}Outside option.} the next best alternative, the thing you fall back on.\par
\textbf{\color{accent}Opportunity cost.} what you give up by taking one action instead of the next best alternative.\par
\textbf{\color{accent}Economic cost.} direct cost plus opportunity cost.\par
\textbf{\color{accent}Economic rent.} the net benefit from the action you chose, minus the opportunity cost.\par
\endgroup


\needspace{12\baselineskip}\section{Would you go to the game?}\setcounter{q}{0}

\Q{Fill in the numbers in the table. Your \textbf{benefit} from the game is the most you would pay for the Dodgers day. The \textbf{cost} of the shift is the least you would accept to work it. Then fill in the net benefit of each option.}

\begin{center}
\renewcommand{\arraystretch}{1.45}
\begin{tabular}{@{}p{1.9in}p{1.35in}p{1.35in}@{}}
\toprule
 & Go to the game & Work the shift \\
\midrule
Benefit & \blank & \$120 \\
Direct cost & \$100 & \blank \\
\midrule
Net benefit & \blank & \blank \\
\bottomrule
\end{tabular}
\end{center}

\Q{Given these numbers, what is your opportunity cost of going to the game, with the shift on the table?}
\answerlines{1}

\Q{What is the economic cost of going to the game?}
\answerlines{1}

\Q{What is the economic cost of taking the shift?}
\answerlines{1}

\Q{So, which option do you choose? What is your economic rent from that option?}
\answerlines{2}

\Q{Now suppose there is no shift on offer: your alternative is a quiet afternoon at home. Do you go to the game now? What is your opportunity cost of going in this case?}
\answerlines{2}

\end{document}
